\chapter{Ensembles et structures}\label{ch:b2:structures}

Ce chapitre d'ouverture affûte les fondations posées dans le volume
de première année pour en faire des outils de travail courants~: le
calcul des ensembles et des quotients, la comparaison des ensembles
infinis (dénombrabilité, Cantor--Bernstein), et la théorie
structurelle des groupes et des anneaux --- le théorème de Lagrange,
le groupe symétrique et sa signature, les idéaux et le théorème des
restes chinois. Tout ici sert sans relâche dans la suite du livre~:
la signature construit le déterminant
(\cref{ch:b2:linalg}), les \omterm{def:b2:structures:quotientring}{anneaux quotients} font marcher
l'arithmétique, et la dénombrabilité sous-tend à la fois la topologie
et les probabilités.

\section{Ensembles, applications, quotients}

Nous utilisons librement le langage des ensembles, des applications
et des relations d'équivalence et d'\omterm{def:b2:structures:generated}{ordre} mis en place dans le volume
de première année. Deux améliorations méritent un énoncé en bonne et
due forme.

\begin{proposition}[Images et images réciproques de familles]\label{prop:b2:structures:images}
Soient $f \colon E \to F$ et $(A_i)_{i \in I}$, $(B_j)_{j \in J}$
des familles de parties de $E$, resp.\ $F$. Alors
\[
f^{-1}\Bigl(\bigcup_j B_j\Bigr) = \bigcup_j f^{-1}(B_j),
\qquad
f^{-1}\Bigl(\bigcap_j B_j\Bigr) = \bigcap_j f^{-1}(B_j),
\qquad
f^{-1}(F \setminus B) = E \setminus f^{-1}(B),
\]
\[
f\Bigl(\bigcup_i A_i\Bigr) = \bigcup_i f(A_i),
\qquad
f\Bigl(\bigcap_i A_i\Bigr) \subseteq \bigcap_i f(A_i)
\quad (\text{égalité pour } f \text{ injective}).
\]
\end{proposition}

\begin{proof}
Chaque identité s'obtient en déroulant les définitions~; par exemple
$x \in f^{-1}(\bigcap B_j) \iff f(x) \in B_j$ pour tout $j$ $\iff x \in
f^{-1}(B_j)$ pour tout $j$. Les identités sur les images et l'échec de
l'égalité dans le cas de l'intersection (avec la correction par
injectivité) ont été démontrés dans le volume de première année pour
deux ensembles~; les arguments sont identiques pour des familles.
\end{proof}

\begin{example}[Où l'inclusion des images est stricte]\label{ex:b2:structures:imagestrict}
Prenons $f \colon \R \to \R$, $f(x) = x^2$, avec $A_1 =
\intcc{-1}{0}$ et $A_2 = \intcc{0}{1}$. Alors
\[
f(A_1 \cap A_2) = f(\{0\}) = \{0\},
\qquad
f(A_1) \cap f(A_2) = \intcc{0}{1} \cap \intcc{0}{1} =
\intcc{0}{1} :
\]
l'inclusion de la~\cref{prop:b2:structures:images} est aussi
stricte que possible --- les deux antécédents $\pm x$ d'une
valeur commune vivent dans des $A_i$ différents. L'injectivité
est précisément ce qui interdit cette séparation, et c'est
pourquoi les images réciproques (qui ne fusionnent jamais de
points) vérifient les quatre identités sans condition, tandis que
les images perdent celle sur les intersections. Règle empirique
pour tout le livre~: faites passer les \emph{images réciproques} à
travers les opérations ensemblistes librement~; maniez les images
avec précaution.
\end{example}

\begin{definition}[Ensemble quotient]\label{def:b2:structures:quotient}
Soit $\mathcal{R}$ une relation d'équivalence sur $E$. L'\emph{ensemble
quotient}\index{ensemble quotient} $E/\mathcal{R}$ est l'ensemble des
classes d'équivalence~; la surjection $\pi \colon E \to E/\mathcal{R}$,
$x \mapsto \mathrm{cl}(x)$, est la \emph{projection
canonique}\index{projection canonique}.

\emph{Propriété universelle (factorisation)~:} si $f \colon E \to F$ est
\emph{compatible} avec $\mathcal{R}$ (c.-à-d.\ $x \mathbin{\mathcal{R}}
y \implies f(x) = f(y)$), il existe une unique application $\overline f
\colon E/\mathcal{R} \to F$ telle que $f = \overline f \circ \pi$.
\end{definition}

\begin{proof}[Démonstration de la propriété universelle]
Unicité~: la condition $f = \overline f \circ \pi$ s'écrit
\[
\overline f\bigl(\mathrm{cl}(x)\bigr) = f(x)
\qquad (x \in E),
\]
et comme $\pi$ est surjective, tout élément de
$E/\mathcal{R}$ est un certain $\mathrm{cl}(x)$~: les valeurs de
$\overline f$ sont toutes imposées. Existence~: on prend la
formule affichée comme \emph{définition} de $\overline f$~; elle
est sans ambiguïté précisément grâce à la compatibilité --- si
$\mathrm{cl}(x) = \mathrm{cl}(y)$, alors $x \mathbin{\mathcal{R}}
y$, donc $f(x) = f(y)$ et les deux valeurs candidates coïncident
--- et elle factorise $f$ par construction. Noter le partage des
tâches~: la surjectivité de $\pi$ donne l'unicité, la
compatibilité donne l'existence.
\end{proof}

\begin{example}
$\Z/n\Z$ est le quotient de $\Z$ par la congruence modulo $n$~; les
vérifications de bonne définition du volume de première année étaient
des cas particuliers de la propriété universelle. Les quotients
transforment les «~constructions compatibles sur les représentants~»
en applications véritables --- nous nous en servons constamment
ci-dessous.
\end{example}

\section{Dénombrabilité et cardinalité}

\begin{definition}[Équipotence, dénombrabilité]\label{def:b2:structures:countable}
Deux ensembles sont \emph{équipotents}\index{équipotence} lorsqu'une
bijection les relie. Un ensemble est \emph{dénombrable}\index{ensemble
dénombrable} lorsqu'il est équipotent à $\N$ (certains auteurs incluent
les ensembles finis~; nous disons \emph{au plus dénombrable} pour
«~fini ou dénombrable~»).
\end{definition}

\begin{proposition}[Propriétés de stabilité]\label{prop:b2:structures:countablestable}
\begin{enumerate}
  \item Toute partie infinie de $\N$ est \omterm{def:b2:structures:countable}{dénombrable}~; un ensemble est
        au plus \omterm{def:b2:structures:countable}{dénombrable} si et seulement s'il s'injecte dans $\N$ si
        et seulement s'il est vide ou image surjective de $\N$.
  \item $\N \times \N$ est \omterm{def:b2:structures:countable}{dénombrable}~; un produit de deux ensembles
        au plus \omterm{def:b2:structures:countable}{dénombrables} est au plus \omterm{def:b2:structures:countable}{dénombrable}.
  \item Une union au plus \omterm{def:b2:structures:countable}{dénombrable} d'ensembles au plus \omterm{def:b2:structures:countable}{dénombrables}
        est au plus \omterm{def:b2:structures:countable}{dénombrable}.
  \item $\Z$ et $\Q$ sont \omterm{def:b2:structures:countable}{dénombrables}.
\end{enumerate}
\end{proposition}

\begin{proof}
(1) Énumérons une partie infinie $A \subseteq \N$ par minima
successifs~: $a_0 = \min A$, $a_{k+1} = \min\,(A \setminus \{a_0, \dots,
a_k\})$ (non vide car $A$ est infini)~; l'application $k \mapsto a_k$
est strictement croissante, injective, et surjective sur $A$ (tout $a
\in A$ ne dépasse qu'un nombre fini d'éléments de $A$, donc il est
atteint). Si $E$ s'injecte dans $\N$ via $\varphi$, alors $E$ est
équipotent à $\varphi(E) \subseteq \N$~: fini ou \omterm{def:b2:structures:countable}{dénombrable}. Si $s
\colon \N \to E$ est surjective, alors $x \mapsto \min s^{-1}(\{x\})$
injecte $E$ dans $\N$.

(2) L'application $(p, q) \mapsto 2^p(2q + 1) - 1$ est une bijection
$\N^2 \to \N$ (tout entier strictement positif possède une unique
décomposition pair--impair $2^p m$ avec $m$ impair, par unicité de la
factorisation). Produits~: composer les injections.

(3) Étant donné des ensembles $E_n$ munis de surjections $s_n \colon \N
\to E_n$ (sans dommage lorsque l'un des $E_n$ est fini~: on répète des
valeurs), l'application $(n, k) \mapsto s_n(k)$ est une surjection du
\omterm{def:b2:structures:countable}{dénombrable} $\N^2$ sur $\bigcup E_n$.

(4) $\Z = \N \cup (-\N^*)$~: union \omterm{def:b2:structures:countable}{dénombrable}. $\Q$ est image
surjective de $\Z \times \N^*$ (l'application fraction), donc au plus
\omterm{def:b2:structures:countable}{dénombrable}, et infini.
\end{proof}

\begin{example}[Une fonction de couplage, calculée]\label{ex:b2:structures:pairing}
La bijection $(p, q) \mapsto 2^p(2q + 1) - 1$ de la démonstration
mérite d'être vue à l'œuvre. Ses premières valeurs~:
\[
\begin{array}{c|ccccc}
 & q = 0 & q = 1 & q = 2 & q = 3 & q = 4\\
\hline
p = 0 & 0 & 2 & 4 & 6 & 8\\
p = 1 & 1 & 5 & 9 & 13 & 17\\
p = 2 & 3 & 11 & 19 & 27 & 35\\
p = 3 & 7 & 23 & 39 & 55 & 71
\end{array}
\]
La ligne $p$ regroupe les entiers $n$ pour lesquels $n + 1$ est
exactement divisible par $2^p$~: tout entier naturel apparaît
exactement une fois. Le décodage est aussi explicite que le
codage~: pour $n = 43$, on factorise $n + 1
= 44 = 2^2\cdot 11 = 2^2(2\cdot5 + 1)$, d'où $(p, q) = (2, 5)$.
L'idée à retenir~: les démonstrations de dénombrabilité sont
souvent des \emph{algorithmes} déguisés --- ici, «~mettre les
puissances de deux en facteur~».
\end{example}

\begin{example}[Les nombres algébriques sont dénombrables]\label{ex:b2:structures:algebraic}
Un nombre complexe est \emph{algébrique} lorsqu'il annule un
polynôme non nul à coefficients rationnels. L'ensemble
$\overline\Q$ des nombres algébriques est \omterm{def:b2:structures:countable}{dénombrable}~: les
polynômes de degré $\leq d$ sur $\Q$ s'injectent dans $\Q^{d+1}$,
un produit fini d'\omterm{def:b2:structures:countable}{ensembles dénombrables}
(\cref{prop:b2:structures:countablestable}~(2))~; l'union sur
$d$ énumère les polynômes rationnels non nuls en $P_0, P_1,
P_2, \dots$~; chaque $P_k$ a un nombre fini de racines~; et
\[
\overline\Q = \bigcup_{k \in \N}\ \{\text{racines de } P_k\}
\]
est une union \omterm{def:b2:structures:countable}{dénombrable} d'ensembles finis
(\cref{prop:b2:structures:countablestable}~(3)), infinie car
elle contient $\Q$. Combiné avec la non-dénombrabilité de $\R$
(\cref{thm:b2:structures:cantor} ci-dessous), cela prouve --- sans
en exhiber un seul --- que les nombres transcendants existent
et forment une majorité non \omterm{def:b2:structures:countable}{dénombrable}~: l'argument de comptage
de Cantor de 1874, l'existence par la seule cardinalité.
\end{example}

\begin{theorem}[Cantor~; non-dénombrabilité de $\R$]\label{thm:b2:structures:cantor}
\begin{enumerate}
  \item Pour tout ensemble $E$, il n'existe pas de surjection $E \to
        \mathcal{P}(E)$.
  \item $\R$ n'est \emph{pas} \omterm{def:b2:structures:countable}{dénombrable}.
\end{enumerate}
\end{theorem}

\begin{proof}
(1) a été démontré dans le volume de première année (l'ensemble
diagonal $D = \{x : x \notin f(x)\}$).

(2) Supposons que $(x_n)_{n \in \N}$ énumère $\R$. Construisons des
segments emboîtés $I_0 \supseteq I_1 \supseteq \dots$ avec $\abs{I_n}
= 3^{-n}$ et $x_n \notin I_n$~: on découpe le segment courant en trois
tiers fermés~; au moins un tiers évite $x_n$ (un point rencontre au
plus deux des trois). Le théorème des segments emboîtés (bornes
adjacentes) fournit $\ell \in \bigcap_n I_n$~; mais $\ell = x_N$ pour
un certain $N$, et $x_N \notin I_N$~: contradiction.
\end{proof}

\begin{theorem}[Cantor--Bernstein]\label{thm:b2:structures:cantorbernstein}
Si $E$ s'injecte dans $F$ et $F$ s'injecte dans $E$, alors $E$ et $F$
sont \omterm{def:b2:structures:countable}{équipotents}.\index{théorème de Cantor--Bernstein}
\end{theorem}

\begin{proof}
Soient $f \colon E \to F$ et $g \colon F \to E$ des injections. Pour
chaque point (de $E$ ou de $F$), suivons sa \emph{chaîne d'ancêtres}
d'images réciproques successives, $x \mapsto g^{-1}(x) \mapsto
f^{-1}(g^{-1}(x)) \mapsto \dots$ --- chaque étape est définie tant que
le point courant appartient à l'image de l'injection concernée, et
est alors unique par injectivité. Trois sorts mutuellement exclusifs~:
la chaîne s'arrête en un point de $E \setminus g(F)$ (\emph{origine
dans $E$}), s'arrête en un point de $F \setminus f(E)$ (\emph{origine
dans $F$}), ou ne s'arrête jamais. Cela partitionne $E = E_E \cup E_F
\cup E_\infty$ et $F = F_E \cup F_F \cup F_\infty$ selon l'origine.

Observons maintenant~: $f$ envoie $E_E$ \emph{sur} $F_E$ --- la chaîne
de $f(x)$ est la chaîne de $x$ précédée d'une étape, donc les origines
coïncident~; et tout $y \in F_E$ possède une chaîne d'au moins une
étape (son origine est dans $E$), donc $y = f(x)$ avec $x \in E_E$. Le
même argument donne des bijections $f \colon E_\infty \to F_\infty$ et
$g \colon F_F \to E_F$. En recollant,
\[
h(x) =
\begin{cases}
f(x) & \text{si } x \in E_E \cup E_\infty,\\
g^{-1}(x) & \text{si } x \in E_F,
\end{cases}
\]
est une bijection de $E$ sur $F = F_E \cup F_\infty \cup F_F$~: elle
est bijective par morceaux, et les trois morceaux d'arrivée sont
disjoints.
\end{proof}

\begin{example}\label{ex:b2:structures:cbexample}
$\intoo{0}{1}$ et $\intcc{0}{1}$ sont \omterm{def:b2:structures:countable}{équipotents}~: l'identité
injecte dans un sens, $x \mapsto \frac{x + 1}{3}$ dans l'autre~; le
théorème fabrique la bijection (nécessairement discontinue). De même
$\R$, $\intoo{0}{1}$ (via des bijections de type $\tanh$) et
$\mathcal{P}(\N)$ (développements binaires, \cref{exo:b2:structures:3})
sont tous \omterm{def:b2:structures:countable}{équipotents}~: «~la puissance du continu~».
\end{example}

\begin{example}[Le segment et le carré]\label{ex:b2:structures:squareseg}
$\intcc{0}{1}$ et $\intcc{0}{1}^2$ sont \omterm{def:b2:structures:countable}{équipotents} --- la dimension
est invisible à la cardinalité. Une injection est triviale~: $x
\mapsto (x, 0)$. Pour l'autre, on envoie $(x, y)$ sur le réel dont
les chiffres décimaux entrelacent ceux de $x$ et $y$,
\[
(0.x_1x_2x_3\dots,\ 0.y_1y_2y_3\dots)
\;\longmapsto\; 0.x_1y_1x_2y_2x_3y_3\dots,
\]
en choisissant pour chaque coordonnée le développement qui ne se
termine pas par une suite infinie de $9$~: avec cette convention,
les chiffres de l'image déterminent ceux de $x$ et de $y$, donc
l'application est injective (elle n'a pas besoin d'être surjective
--- les images n'ont, par exemple, jamais leurs chiffres de
position impaire égaux à $9$ à partir d'un certain rang --- et cela
ne pose pas de problème). Le théorème de Cantor--Bernstein
(\cref{thm:b2:structures:cantorbernstein}) assemble une véritable
bijection. La continuité, bien sûr, est sans espoir~: une bijection
continue entre les deux est impossible --- les chapitres de
topologie métrique expliquent pourquoi (la connexité distingue la
droite du plan, \cref{ch:b2:metric}).
\end{example}

\section{Groupes}

\begin{definition}[Sous-groupe engendré~; ordre]\label{def:b2:structures:generated}
Soit $G$ un groupe et $A \subseteq G$. Le sous-groupe \emph{engendré}
par $A$, noté $\langle A \rangle$, est le plus petit sous-groupe
contenant $A$ --- concrètement, tous les produits finis d'éléments de
$A$ et de leurs inverses. Un groupe est \emph{cyclique}\index{groupe
cyclique} lorsqu'il est engendré par un seul élément~: $\langle a\rangle
= \{a^k : k \in \Z\}$. L'\emph{ordre}\index{ordre!d'un élément} de $a
\in G$ est $\operatorname{ord}(a) = \abs{\langle a \rangle}$
(éventuellement infini)~; lorsqu'il est fini, c'est le plus petit $n
\geq 1$ tel que $a^n = e$, et $a^k = e \iff \operatorname{ord}(a) \mid
k$.
\end{definition}

\begin{proof}[Démonstration de la caractérisation de l'ordre]
S'il existe $a^m = e$ avec $m \geq 1$, soit $n \geq 1$ le plus petit
tel que $a^n = e$. Les éléments $e, a, \dots, a^{n-1}$ sont deux à deux
distincts ($a^{i} = a^{j}$ avec $0 \leq i < j < n$ donne $a^{j-i} = e$,
contredisant la minimalité), et tout $a^k$ se ramène à l'un d'eux par
division euclidienne $k = nq + r$~: $\langle a\rangle$ a exactement $n$
éléments, et $a^k = a^r = e \iff r = 0 \iff n \mid k$. Si aucune
puissance n'est triviale, tous les $a^k$ ($k \in \Z$) sont distincts
(même argument de division) et l'\omterm{def:b2:structures:generated}{ordre} est infini.
\end{proof}

\begin{theorem}[Lagrange]\label{thm:b2:structures:lagrange}
Soit $G$ un groupe fini et $H$ un sous-groupe. Alors $\abs H$ divise
$\abs G$.\index{théorème de Lagrange} En particulier, l'\omterm{def:b2:structures:generated}{ordre} de tout
élément divise $\abs G$, et $a^{\abs G} = e$ pour tout $a \in G$.
\end{theorem}

\begin{proof}
La relation $x \sim y \iff x^{-1}y \in H$ est une équivalence
(réflexive~: $e \in H$~; symétrique~: par les inverses~; transitive~:
par les produits). La classe de $x$ est la \emph{classe à gauche} $xH
= \{xh : h \in H\}$, et $h \mapsto xh$ est une bijection $H \to xH$
(inverse $y \mapsto x^{-1}y$)~: toutes les classes ont $\abs H$
éléments. Les classes partitionnent $G$ (le théorème général de
partition du volume de première année), donc $\abs G = \abs H \times
(\text{nombre de classes})$. Pour un élément~: on applique ceci à $H =
\langle a\rangle$~; alors $a^{\abs G} = (a^{\operatorname{ord}
a})^{\abs G / \operatorname{ord} a} = e$.
\end{proof}

\begin{example}[Les classes à l'œuvre~: $A_3$ dans $\mathfrak{S}_3$]\label{ex:b2:structures:cosets}
Prenons $G = \mathfrak{S}_3$ (\omterm{def:b2:structures:generated}{ordre} $6$) et $H = A_3 =
\{\mathrm{id},\ (1\,2\,3),\ (1\,3\,2)\}$. Les classes à gauche sont
\[
H = \{\mathrm{id},\ (1\,2\,3),\ (1\,3\,2)\},
\qquad
(1\,2)H = \{(1\,2),\ (2\,3),\ (1\,3)\} :
\]
deux classes de trois éléments partitionnant $G$, exactement comme
l'exige le décompte $\abs G = \abs H \times (\text{nombre de
classes})$ --- et visiblement la partition en permutations paires
et impaires. Noter que $(1\,3)H = (1\,2)H$ bien que $(1\,3) \neq
(1\,2)$~: les classes sont des \emph{classes}, non repérées par
leurs représentants, et $x^{-1}y \in H$ est la seule comparaison
légitime. Cette image à deux classes est celle qui vaut en général
pour la signature~: $A_n$ et son unique classe compagne coupent
$\mathfrak{S}_n$ en deux, ce qui est la façon dont le problème du
week-end compte les positions du puzzle atteignables.
\end{example}

\begin{example}\label{ex:b2:structures:lagrangeapp}
Deux dividendes immédiats. \emph{Les groupes d'\omterm{def:b2:structures:generated}{ordre} premier sont
\omterm{def:b2:structures:generated}{cycliques}~:} si $\abs G = p$ est premier et $a \neq e$, alors
$\operatorname{ord}(a)$ divise $p$ et n'est pas $1$, donc vaut $p$~:
$\langle a\rangle = G$. \emph{Le treillis des sous-groupes de
$\Z/12\Z$~:} d'après la~\cref{prop:b2:structures:cyclic} ci-dessous,
il y a exactement un sous-groupe par diviseur de $12$ --- d'\omterm{def:b2:structures:generated}{ordres}
$1, 2, 3, 4, 6, 12$, engendrés respectivement par $\overline 0$,
$\overline 6$, $\overline 4$, $\overline 3$, $\overline 2$,
$\overline 1$. La mise en garde finale~: la \emph{réciproque} de
Lagrange est fausse en général --- $A_4$ est d'\omterm{def:b2:structures:generated}{ordre} $12$ mais n'a
pas de sous-groupe d'\omterm{def:b2:structures:generated}{ordre} $6$, comme nous le démontrons dans le
problème du week-end de ce chapitre (\cref{pb:b2:structures:1},
question 14). Lagrange restreint les \omterm{def:b2:structures:generated}{ordres} possibles~; il ne les
garantit pas.
\end{example}

\begin{omfigure}
\begin{tikzpicture}[scale=1.05, every node/.style={font=\small}]
  \node (G) at (0,3)  {$\Z/12\Z = \langle\overline1\rangle$};
  \node (A) at (-1.6,2) {$\langle\overline2\rangle$};
  \node (B) at (1.6,2)  {$\langle\overline3\rangle$};
  \node (C) at (-1.6,1) {$\langle\overline4\rangle$};
  \node (D) at (1.6,1)  {$\langle\overline6\rangle$};
  \node (E) at (0,0)  {$\{\overline0\}$};
  \draw[omDef] (G) -- (A); \draw[omDef] (G) -- (B);
  \draw[omDef] (A) -- (C); \draw[omDef] (A) -- (D);
  \draw[omDef] (B) -- (D); \draw[omDef] (C) -- (E);
  \draw[omDef] (D) -- (E);
  \node[omThm, right=2pt] at (0.1,3) {\ ordre $12$};
  \node[omThm, left=2pt]  at (-2.1,2) {ordre $6$\ };
  \node[omThm, right=2pt] at (2.1,2) {\ ordre $4$};
  \node[omThm, left=2pt]  at (-2.1,1) {ordre $3$\ };
  \node[omThm, right=2pt] at (2.1,1) {\ ordre $2$};
\end{tikzpicture}

{\small Le treillis des sous-groupes de $\Z/12\Z$~: un sous-groupe
par diviseur de $12$ (\cref{prop:b2:structures:cyclic}), avec une
arête lorsque l'un contient l'autre avec un indice premier. Les
inclusions vont \emph{à rebours} de la divisibilité du générateur~:
$\langle\overline 4\rangle \subseteq \langle\overline2\rangle$ car
$4$ est un multiple de $2$.}
\end{omfigure}

\begin{proposition}[Groupes cycliques]\label{prop:b2:structures:cyclic}
Soit $G = \langle a \rangle$ \omterm{def:b2:structures:generated}{cyclique} d'\omterm{def:b2:structures:generated}{ordre} $n$.
\begin{enumerate}
  \item $G$ est isomorphe à $(\Z/n\Z, +)$, via $\overline k \mapsto
        a^k$.
  \item Tout sous-groupe de $G$ est \omterm{def:b2:structures:generated}{cyclique}~; pour chaque diviseur $d
        \mid n$ il existe exactement un sous-groupe d'\omterm{def:b2:structures:generated}{ordre} $d$, à
        savoir $\langle a^{n/d}\rangle$.
  \item $a^k$ engendre $G$ si et seulement si $\gcd(k, n) = 1$~: $G$ a
        $\varphi(n)$ générateurs (indicatrice d'Euler\index{indicatrice
        d'Euler}).
\end{enumerate}
\end{proposition}

\begin{proof}
(1) L'application $k \mapsto a^k$ de $\Z$ sur $G$ est compatible avec
la congruence modulo $n$ ($a^{k} = a^{k'} \iff n \mid k - k'$, par la
caractérisation de l'\omterm{def:b2:structures:generated}{ordre})~; la propriété universelle
(\cref{def:b2:structures:quotient}) fournit un morphisme bijectif bien
défini depuis $\Z/n\Z$.

(2) Soit $H \leq G$ non trivial et $m$ le plus petit $\geq 1$ tel que
$a^m \in H$. La division euclidienne montre $H = \langle a^m\rangle$
(pour $a^k \in H$~: $k = mq + r$ force $a^r \in H$, donc $r = 0$), et
$m \mid n$ (on divise $n$ par $m$~: $a^{n \bmod m} \in H$). Alors
$\abs H = n/m$~; prendre $m = n/d$ réalise chaque diviseur $d$.
Unicité~: tout sous-groupe d'\omterm{def:b2:structures:generated}{ordre} $d$ est, d'après ce qui précède,
de la forme $\langle a^m \rangle$ avec $n/m = d$ --- donc $m = n/d$
est imposé et le sous-groupe est déterminé.

(3) Nous affirmons que $\operatorname{ord}(a^k) = \frac{n}{\gcd(k,
n)}$. Posons $d = \gcd(k, n)$. Pour tout $m \geq 1$, la
caractérisation de l'\omterm{def:b2:structures:generated}{ordre} de la~\cref{def:b2:structures:generated}
donne la chaîne d'équivalences
\[
(a^k)^m = e
\iff n \mid km
\iff \frac{n}{d} \,\Big|\, \frac{k}{d}\,m
\iff \frac{n}{d} \,\Big|\, m ,
\]
la dernière étape par le lemme de Gauss, puisque $\frac nd$ et
$\frac kd$ sont premiers entre eux. Le plus petit tel $m$ est
$\frac nd$~: $\operatorname{ord}(a^k) = \frac n{\gcd(k,n)}$, qui
vaut $n$ si et seulement si $\gcd(k, n) = 1$. Il y a $\varphi(n)$
telles classes $k$ modulo $n$.
\end{proof}

\section{Le groupe symétrique}

\begin{definition}\label{def:b2:structures:sn}
$\mathfrak{S}_n$ est le groupe des permutations de $\intint{1}{n}$
(d'\omterm{def:b2:structures:generated}{ordre} $n!$). Un \emph{cycle}\index{cycle} $(a_1\,a_2\,\cdots\,a_k)$
envoie $a_1 \mapsto a_2 \mapsto \dots \mapsto a_k \mapsto a_1$ et
fixe tout le reste~; $k$ est sa \emph{longueur}, un $2$-cycle est une
\emph{transposition}\index{transposition}. Deux cycles sont
\emph{disjoints} lorsque leurs
supports (points non fixes) le sont.
\end{definition}

\begin{theorem}[Décomposition en cycles]\label{thm:b2:structures:cycles}
Toute permutation $\sigma \neq \mathrm{id}$ est un produit de \omterm{def:b2:structures:sn}{cycles}
deux à deux disjoints, de manière unique à l'\omterm{def:b2:structures:generated}{ordre} des facteurs près.
Des \omterm{def:b2:structures:sn}{cycles} disjoints commutent, et $\operatorname{ord}(\sigma)$ est le
ppcm des longueurs.
\end{theorem}

\begin{proof}
Considérons la relation d'«~orbite~» sur le support de $\sigma$~: $x
\sim y$ si et seulement si $y = \sigma^k(x)$ pour un certain $k \in \Z$
--- une relation d'équivalence. Chaque classe $\{x, \sigma(x), \dots,
\sigma^{k-1}(x)\}$ (finie, donc les itérées bouclent --- la première
répétition doit revenir à $x$ par injectivité) porte le \omterm{def:b2:structures:sn}{cycle} $(x\
\sigma(x)\ \cdots\ \sigma^{k-1}(x))$, et $\sigma$ est le produit de ces
\omterm{def:b2:structures:sn}{cycles}~: sur chaque orbite, seul le \omterm{def:b2:structures:sn}{cycle} correspondant agit.
Unicité~: toute factorisation en \omterm{def:b2:structures:sn}{cycles} disjoints reproduit exactement
les orbites (le \omterm{def:b2:structures:sn}{cycle} passant par $x$ doit être $(x\ \sigma(x)\
\cdots)$). Des \omterm{def:b2:structures:sn}{cycles} disjoints commutent puisqu'ils déplacent des
points disjoints~; l'énoncé sur l'\omterm{def:b2:structures:generated}{ordre} s'ensuit car $\sigma^m =
\mathrm{id}$ si et seulement si la puissance $m$-ième de chaque \omterm{def:b2:structures:sn}{cycle}
l'est (disjonction), si et seulement si chaque longueur divise $m$.
\end{proof}

\begin{example}[Le type de cycle comme recensement]\label{ex:b2:structures:cycletype}
Combien de permutations de $\mathfrak{S}_9$ ont le type de \omterm{def:b2:structures:sn}{cycle}
$(4, 3, 2)$ --- un $4$-cycle, un $3$-cycle, une \omterm{def:b2:structures:sn}{transposition}~?
On choisit les supports et les \omterm{def:b2:structures:generated}{ordres} \omterm{def:b2:structures:generated}{cycliques}~:
\[
\frac{9!}{4\cdot 3\cdot 2}
= \frac{362\,880}{24} = 15\,120 :
\]
on aligne les neuf symboles en ligne ($9!$ façons), on regroupe
les quatre premiers, les trois suivants, les deux derniers en
\omterm{def:b2:structures:sn}{cycles}, et on divise par les rotations à l'intérieur de chaque
groupe ($4$, $3$ et $2$ respectivement) qui donnent la même
permutation. (Ici les \emph{longueurs} des \omterm{def:b2:structures:sn}{cycles} sont distinctes,
donc pas d'autre division~; des longueurs égales exigeraient aussi
de diviser par les permutations des groupes de même taille.) Toute
telle permutation est d'\omterm{def:b2:structures:generated}{ordre} $\operatorname{lcm}(4,3,2) = 12$ et
de signature $(-1)^3(-1)^2(-1)^1 = +1$
(\cref{thm:b2:structures:cycles} et le théorème de la signature
ci-dessous). Une partition de $9$, une classe de conjugaison, un
recensement --- la combinatoire de $\mathfrak{S}_n$ est
l'arithmétique des partitions.
\end{example}

\begin{theorem}[Signature]\label{thm:b2:structures:signature}
Il existe exactement un morphisme de groupes $\varepsilon \colon
\mathfrak{S}_n \to \{\pm 1\}$ (pour $n \geq 2$) prenant la valeur $-1$
sur les \omterm{def:b2:structures:sn}{transpositions}~: la \emph{signature}\index{signature}. De plus
$\varepsilon(\sigma) = (-1)^{I(\sigma)}$ où $I(\sigma)$ est le nombre
d'\emph{inversions} (paires $i < j$ avec $\sigma(i) > \sigma(j)$), un
$k$-cycle a pour signature $(-1)^{k-1}$, et le \emph{groupe alterné}
$A_n = \ker\varepsilon$ est d'\omterm{def:b2:structures:generated}{ordre} $\frac{n!}{2}$.
\end{theorem}

\begin{proof}
\emph{Existence.} Pour $\sigma \in \mathfrak{S}_n$ posons
\[
\varepsilon(\sigma)
= \prod_{1 \leq i < j \leq n}
\frac{\sigma(j) - \sigma(i)}{j - i} .
\]
Les valeurs absolues des facteurs se multiplient en $1$ (les paires
non ordonnées $\{\sigma(i), \sigma(j)\}$ parcourent toutes les
paires), donc $\varepsilon(\sigma) = (-1)^{I(\sigma)} \in \{\pm1\}$.
Morphisme~: pour $\sigma, \tau$,
\[
\varepsilon(\sigma\tau)
= \prod_{i<j} \frac{\sigma(\tau(j)) - \sigma(\tau(i))}{j - i}
= \prod_{i<j} \frac{\sigma(\tau(j)) - \sigma(\tau(i))}{\tau(j) -
\tau(i)} \cdot \prod_{i<j} \frac{\tau(j) - \tau(i)}{j - i}
= \varepsilon(\sigma)\,\varepsilon(\tau),
\]
le produit du milieu valant $\varepsilon(\sigma)$ après réindexation
par les paires $\{\tau(i), \tau(j)\}$ (chaque paire non ordonnée
apparaît une fois, et numérateur et dénominateur changent de signe
ensemble). Une \omterm{def:b2:structures:sn}{transposition} $\tau = (a\,b)$ avec $a < b$ a un nombre
impair d'inversions~; comptées exactement~: les paires inversées $(i,
j)$, $i < j$, avec $\tau(i) > \tau(j)$ sont
\[
(a, j) \ \text{pour } a < j < b, \qquad
(i, b) \ \text{pour } a < i < b, \qquad
(a, b) \ \text{elle-même},
\]
soit $(b - a - 1) + (b - a - 1) + 1 = 2(b - a) - 1$ d'entre elles,
un nombre impair. (Autrement~: vérifier $(1\,2)$ directement, avec
une inversion, et conjuguer --- les conjugués ont même signature
puisque $\varepsilon$ est un morphisme vers un groupe abélien.) Donc
$\varepsilon((a\,b)) = (-1)^{2(b-a)-1} = -1$.

\emph{Unicité.} Les \omterm{def:b2:structures:sn}{transpositions} engendrent $\mathfrak{S}_n$ (tout
\omterm{def:b2:structures:sn}{cycle} $(a_1\cdots a_k) = (a_1\,a_k)(a_1\,a_{k-1})\cdots(a_1\,a_2)$,
et le~\cref{thm:b2:structures:cycles} conclut)~; un morphisme vers
$\{\pm1\}$ est déterminé par ses valeurs sur des générateurs.

\emph{Conséquences.} L'identité sur les \omterm{def:b2:structures:sn}{cycles} ci-dessus écrit un
$k$-cycle comme $k - 1$ \omterm{def:b2:structures:sn}{transpositions}~: signature $(-1)^{k-1}$.
$A_n$~: le morphisme $\varepsilon$ est surjectif (des \omterm{def:b2:structures:sn}{transpositions}
existent pour $n \geq 2$), et les deux «~classes~» $A_n$ et
$(1\,2)A_n$ sont équipotentes et partitionnent $\mathfrak{S}_n$
(argument de Lagrange)~: $\abs{A_n} = \frac{n!}{2}$.
\end{proof}

\begin{example}\label{ex:b2:structures:snexample}
$\sigma = \begin{pmatrix} 1&2&3&4&5&6\\ 3&6&5&4&1&2 \end{pmatrix}
= (1\,3\,5)(2\,6)$~: \omterm{def:b2:structures:generated}{ordre} $\operatorname{lcm}(3,2) = 6$, signature
$(-1)^{2}\cdot(-1)^{1} = -1$. La signature est le test de parité le
plus rapide sur les battages --- et le moteur du déterminant au
\cref{ch:b2:linalg}.
\end{example}

\begin{example}[Trois chemins vers un même signe]\label{ex:b2:structures:threeroads}
Soit $\sigma \in \mathfrak{S}_5$ envoyant $1, 2, 3, 4, 5$ sur $3, 5,
4, 1, 2$. \emph{Par les \omterm{def:b2:structures:sn}{cycles}~:} $1 \mapsto 3 \mapsto 4 \mapsto 1$
et $2 \mapsto 5 \mapsto 2$, donc $\sigma = (1\,3\,4)(2\,5)$ et
$\varepsilon(\sigma) = (-1)^{2}(-1)^{1} = -1$. \emph{Par les
inversions~:} dans la liste de valeurs $3, 5, 4, 1, 2$ les paires en
désordre sont $(3,1)$, $(3,2)$, $(5,4)$, $(5,1)$, $(5,2)$, $(4,1)$,
$(4,2)$~: sept d'entre elles, et $(-1)^7 = -1$. \emph{Par les
\omterm{def:b2:structures:sn}{transpositions}~:} $\sigma = (1\,4)(1\,3)(2\,5)$, trois facteurs,
$(-1)^3 = -1$. Trois calculs, une parité~: l'unicité dans
le~\cref{thm:b2:structures:signature} garantit qu'aucun schéma de
comptabilité ne peut jamais les faire diverger --- ce qui est
exactement ce qui rend $\varepsilon$ utilisable comme invariant
(voir le problème du week-end).
\end{example}

\begin{remark}[Ce que devient la signature par la suite]
La signature est la graine de trois récoltes ultérieures~: elle
construit le déterminant et sa règle de produit au~\cref{ch:b2:linalg}~;
elle alimente des invariants de parité pour des casse-têtes
combinatoires (le problème du week-end de ce chapitre résout le
taquin grâce à elle)~; et les groupes alternés $A_n$ qu'elle définit
deviennent centraux dans le volume de troisième année, où leur
simplicité pour $n \geq 5$ explique pourquoi les équations de degré
$5$ n'ont pas de solution par radicaux.
\end{remark}

\section{Anneaux, idéaux, quotients}

\begin{definition}[Idéal]\label{def:b2:structures:ideal}
Soit $A$ un anneau commutatif. Un \emph{idéal}\index{idéal} $I
\subseteq A$ est un sous-groupe additif tel que $a x \in I$ pour tout
$a \in A$, $x \in I$. Les noyaux de morphismes d'anneaux sont des
idéaux~; $I = A$ si et seulement si $1 \in I$ si et seulement si $I$
contient un inversible. L'idéal \emph{engendré} par $x$ est $xA =
\{xa\}$ (un idéal \emph{principal}).
\end{definition}

\begin{theorem}[{Idéaux de $\Z$ et de $K[X]$}]\label{thm:b2:structures:principal}
Tout \omterm{def:b2:structures:ideal}{idéal} de $\Z$ est $n\Z$ pour un unique $n \in \N$~; tout \omterm{def:b2:structures:ideal}{idéal} de
$K[X]$ ($K$ un corps) est $P\,K[X]$ pour un unique $P$ unitaire (ou
nul). Par conséquent, les pgcd existent dans les deux anneaux avec des
relations de Bézout~: $x\Z + y\Z = \gcd(x,y)\Z$, et de même pour les
polynômes.
\end{theorem}

\begin{proof}
Pour $\Z$ c'était le théorème des sous-groupes du volume de première
année (un \omterm{def:b2:structures:ideal}{idéal} est en particulier un sous-groupe, et $n\Z$ est un
\omterm{def:b2:structures:ideal}{idéal}). Pour $K[X]$~: soit $I \neq \{0\}$ un \omterm{def:b2:structures:ideal}{idéal} et $P \in I$ non nul
de degré minimal, normalisé unitaire. Pour $F \in I$, la division
euclidienne $F = PQ + R$ donne $R = F - PQ \in I$ avec $\deg R < \deg
P$~: la minimalité force $R = 0$, donc $I = P\,K[X]$. Unicité~: deux
générateurs unitaires se divisent mutuellement. Les énoncés de Bézout
sont l'égalité de l'\omterm{def:b2:structures:ideal}{idéal} $x\Z + y\Z$ (resp.\ son analogue polynomial)
avec l'\omterm{def:b2:structures:ideal}{idéal} principal du pgcd --- la définition même du pgcd utilisée
en première année, reconnue maintenant comme un énoncé sur les idéaux.
\end{proof}

\begin{example}[Un pgcd de polynômes, de deux façons]\label{ex:b2:structures:polygcd}
Calculons $\gcd(X^3 - 1,\ X^2 - 1)$ dans $\Q[X]$. \emph{Par
Euclide~:}
\[
X^3 - 1 = X\,(X^2 - 1) + (X - 1),
\qquad
X^2 - 1 = (X + 1)(X - 1) + 0 ,
\]
donc le pgcd est $X - 1$, et la remontée donne la relation de
Bézout
\[
X - 1 = 1\cdot(X^3 - 1) - X\cdot(X^2 - 1).
\]
\emph{Par les idéaux~:} l'\omterm{def:b2:structures:ideal}{idéal} $(X^3 - 1)\Q[X] + (X^2 - 1)\Q[X]$
est principal (\cref{thm:b2:structures:principal})~; il contient $X
- 1$ (la formule affichée) et est contenu dans $(X - 1)\Q[X]$ (les
deux générateurs s'annulent en $1$, donc sont multiples de $X -
1$)~: le générateur unitaire est $X - 1$. L'idée à retenir~: le
point de vue des idéaux identifie le pgcd \emph{sans diviser} ---
les racines communes localisent l'\omterm{def:b2:structures:ideal}{idéal}, et Euclide ne fait que le
certifier.
\end{example}

\begin{definition}[Anneau quotient $\Z/n\Z$, revisité]\label{def:b2:structures:quotientring}
Pour un \omterm{def:b2:structures:ideal}{idéal} $I$ de $A$, la relation $x \sim y \iff x - y \in I$ est
une équivalence compatible avec $+$ et $\times$~; l'\omterm{def:b2:structures:quotient}{ensemble quotient}
$A/I$ hérite d'une structure d'anneau --- l'\emph{anneau
quotient}\index{anneau quotient} --- faisant de $\pi \colon A \to A/I$
un morphisme de noyau $I$. Pour $A = \Z$, $I = n\Z$ c'est le
$\Z/n\Z$ du volume de première année, désormais muni de sa propriété
universelle~: tout morphisme annulant $I$ se factorise à travers $A/I$.
\end{definition}

\begin{theorem}[Théorème des restes chinois, forme anneau]\label{thm:b2:structures:crt}
Si $\gcd(m, n) = 1$, l'application
\[
\Z/mn\Z \longrightarrow \Z/m\Z \times \Z/n\Z,
\qquad
\overline{x} \longmapsto (x \bmod m,\; x \bmod n)
\]
est un isomorphisme d'anneaux.\index{théorème des restes chinois} Par
conséquent $\varphi(mn) = \varphi(m)\varphi(n)$ pour $m, n$ premiers
entre eux, et
\[
\varphi(n) = n \prod_{p \mid n} \Bigl(1 - \frac 1p\Bigr)
\quad (p \text{ premier}).
\]
\end{theorem}

\begin{proof}
L'application est un morphisme d'anneaux bien défini (les
compatibilités sont immédiates). Injectivité~: $x \equiv 0$ modulo $m$
et modulo $n$ avec $\gcd(m,n) = 1$ force $mn \mid x$ (Gauss).
Surjectivité~: les deux membres ont $mn$ éléments, donc l'injectivité
suffit (cardinaux finis égaux) --- ou explicitement~: à partir d'une
relation de Bézout $um + vn = 1$, la classe de
\[
x = b\,um + a\,vn
\]
s'envoie sur $(a \bmod m,\ b \bmod n)$, puisque $vn = 1 - um \equiv 1
\pmod m$ fait que $x \equiv a \pmod m$, et symétriquement modulo $n$
--- la recette utilisée numériquement dans
l'\cref{ex:b2:structures:crtinverse}. Les inversibles correspondent
aux couples d'inversibles (les inversibles d'un anneau produit sont
les couples d'inversibles), donc $\varphi(mn) =
\varphi(m)\varphi(n)$. Pour une puissance de premier, $\varphi(p^k) =
p^k - p^{k-1}$ (les non-inversibles modulo $p^k$ sont les multiples de
$p$)~; la multiplicativité assemble la formule du produit.
\end{proof}

\begin{example}[Inverser l'isomorphisme chinois]\label{ex:b2:structures:crtinverse}
Prenons $m = 8$, $n = 9$. L'inverse de l'isomorphisme est rendu
explicite par les deux \emph{idempotents}~: on cherche $u \equiv 1
\pmod 8$, $u \equiv 0 \pmod 9$ et $v \equiv 0 \pmod 8$, $v \equiv 1
\pmod 9$. De $u = 9k \equiv 1 \pmod 8$~: $k \equiv 1$, donc $u = 9$~;
de $v = 8k \equiv 1 \pmod 9$~: $-k \equiv 1$, $k \equiv 8$, donc $v =
64$. Alors la classe de $x = 9a + 64b$ modulo $72$ est l'unique
solution de $x \equiv a \pmod 8$, $x \equiv b \pmod 9$~: pour $a =
3$, $b = 5$ on obtient $27 + 320 = 347 \equiv 59 \pmod{72}$ ---
exactement la valeur intermédiaire trouvée par substitution dans
l'\cref{exo:b2:structures:8}. L'idée à retenir~: $u$ et $v$
vérifient $u + v \equiv 1$, $uv \equiv 0$, $u^2 \equiv u$, $v^2
\equiv v$ modulo $72$~; ce sont les images de $(1, 0)$ et $(0,
1)$, et toute décomposition chinoise est au fond une décomposition
de $1$ en idempotents orthogonaux.
\end{example}

\begin{theorem}[Euler~; Fermat revisité]\label{thm:b2:structures:euler}
Les inversibles de $\Z/n\Z$ forment un groupe d'\omterm{def:b2:structures:generated}{ordre} $\varphi(n)$~;
d'où pour $\gcd(a, n) = 1$~:
\[
a^{\varphi(n)} \equiv 1 \pmod n
\qquad (\text{théorème d'Euler}\index{théorème d'Euler}),
\]
et le petit théorème de Fermat est le cas $n = p$ premier, désormais à
une ligne de Lagrange.
\end{theorem}

\begin{proof}
Les classes inversibles sont exactement celles des entiers premiers à
$n$ (volume de première année)~: au nombre de $\varphi(n)$, formant un
groupe pour la multiplication. Lagrange
(\cref{thm:b2:structures:lagrange})~: tout élément à la puissance de
l'\omterm{def:b2:structures:generated}{ordre} du groupe est l'identité.
\end{proof}

\begin{example}[Un groupe d'inversibles sans générateur]\label{ex:b2:structures:units15}
Le groupe $(\Z/15\Z)^*$ a $\varphi(15) = \varphi(3)\varphi(5)
= 8$ éléments. Est-il \omterm{def:b2:structures:generated}{cyclique}~? Calculons les \omterm{def:b2:structures:generated}{ordres} à l'aide de
l'isomorphisme chinois $(\Z/15\Z)^* \simeq (\Z/3\Z)^* \times
(\Z/5\Z)^*$ (un inversible modulo $15$ est un couple
d'inversibles)~: les facteurs ont pour \omterm{def:b2:structures:generated}{ordres} $2$ et $4$, donc
l'\omterm{def:b2:structures:generated}{ordre} de tout élément divise $\operatorname{lcm}(2, 4) = 4 < 8$
--- aucun élément n'engendre. Concrètement~:
\[
2^4 = 16 \equiv 1, \qquad
4^2 = 16 \equiv 1, \qquad
7^4 \equiv 1, \qquad
11^2 = 121 \equiv 1, \qquad
14^2 \equiv 1 \pmod{15} :
\]
des \omterm{def:b2:structures:generated}{ordres} $4, 2, 4, 2, 2$ et jamais $8$. À comparer avec
l'\cref{exo:b2:structures:10}~: $(\Z/p\Z)^*$ \emph{est} \omterm{def:b2:structures:generated}{cyclique}
pour $p$ premier, car là le groupe des inversibles siège dans un
corps. Le théorème d'Euler s'applique toujours avec l'exposant
$\varphi(15) = 8$, mais le véritable exposant universel ici est $4$
--- Euler est une borne supérieure, pas toujours la plus fine.
\end{example}

\begin{definition}[Algèbre]\label{def:b2:structures:algebra}
Une \emph{$K$-algèbre}\index{algèbre} est un $K$-espace vectoriel $A$
muni d'une structure d'anneau dont la multiplication est
$K$-bilinéaire. Exemples~: $K[X]$, $\mathcal{M}_n(K)$, $\mathcal{L}(E)$,
les espaces de fonctions $\mathcal{F}(X, K)$, $\C$ comme $\R$-algèbre.
Les morphismes d'algèbres sont les morphismes d'anneaux linéaires~;
l'\emph{évaluation} $P \mapsto P(u)$ de $K[X]$ vers $\mathcal{L}(E)$
(ou $\mathcal{M}_n(K)$) est l'exemple central, moteur du
\cref{ch:b2:reduction}.
\end{definition}

\begin{example}[Un morphisme d'évaluation et son noyau]\label{ex:b2:structures:evaluation}
Prenons $A = \begin{pmatrix}0 & 1\\ 0 & 0\end{pmatrix}$ et
l'évaluation $\varepsilon_A \colon \R[X] \to \mathcal{M}_2(\R)$,
$P \mapsto P(A)$. Puisque $A^2 = 0$,
\[
P(A) = P(0)\,I + P'(0)\,A =
\begin{pmatrix} P(0) & P'(0)\\ 0 & P(0)\end{pmatrix},
\]
(seuls les termes constant et linéaire de $P$ subsistent). Donc
$\ker\varepsilon_A = \{P : P(0) = P'(0) = 0\} = X^2\,\R[X]$~: un
\omterm{def:b2:structures:ideal}{idéal} principal, exactement comme le prévoit
le~\cref{thm:b2:structures:principal}, engendré par le polynôme
unitaire $X^2$ de plus petit degré dans le noyau --- le
\emph{polynôme minimal} de $A$, vedette du~\cref{ch:b2:reduction}.
L'image est l'\omterm{def:b2:structures:algebra}{algèbre} commutative de dimension deux $\{aI + bA\}$~:
les morphismes d'évaluation réduisent le $\R[X]$ de dimension
infinie à de petites \omterm{def:b2:structures:algebra}{algèbres} calculables.
\end{example}

\begin{remark}[Perspectives~: trois mélodies à écouter]\label{rem:b2:structures:perspectives}
Trois idées structurelles de ce chapitre reviennent tout au long du
volume, chaque fois dans une orchestration plus fournie.
\emph{Factorisation par un quotient}
(\cref{def:b2:structures:quotient})~: elle construit $\Z/n\Z$ ici,
définit des applications sur les ensembles de solutions de systèmes
linéaires au~\cref{ch:b2:linalg}, et sous-tend en silence tout
argument de «~bonne définition sur les classes~». \emph{Invariants}~:
la signature est un morphisme vers $\{\pm1\}$ qu'aucun coup légal ne
peut esquiver --- la même logique donne la règle de produit du
déterminant (\cref{ch:b2:linalg}), l'invariance de la trace par
similitude, et les quantités conservées du~\cref{ch:b2:diffeq}.
\emph{Compter contre une structure}~: Lagrange compte à travers les
classes, la dimension compte à travers les bases
(\cref{ch:b2:linalg}), la multiplicité compte à travers les degrés
des polynômes (\cref{ch:b2:reduction})~; chaque fois qu'une borne
paraît miraculeuse, une partition ou une graduation fait le comptage.
\end{remark}

\begin{remark}[Pièges classiques]\label{rem:b2:structures:pitfalls}
Quatre classiques. (i) Une application sur un quotient doit être
vérifiée \emph{bien définie}~: «~$\overline x \mapsto$ (formule sur
$x$)~» n'est légitime que si la formule est constante sur les
classes --- la compatibilité de la~\cref{def:b2:structures:quotient},
pas une formalité. (ii) $\operatorname{ord}(ab) =
\operatorname{lcm}(\operatorname{ord}a, \operatorname{ord}b)$
est \emph{faux} en général, même pour des éléments qui commutent
($a$ et $a^{-1}$)~; l'\cref{exo:b2:structures:4} donne l'énoncé
correct (premiers entre eux et commutant), et les \omterm{def:b2:structures:sn}{cycles} disjoints
la version correcte pour les permutations. (iii) La dénombrabilité
survit aux \emph{unions} \omterm{def:b2:structures:countable}{dénombrables} et aux \emph{produits} finis,
mais pas aux produits \omterm{def:b2:structures:countable}{dénombrables}~: $\{0,1\}^{\N}$ est non
\omterm{def:b2:structures:countable}{dénombrable} (\cref{exo:b2:structures:3}) bien que chaque facteur ait
deux éléments. (iv) Cantor--Bernstein ne réclame que des injections
dans les deux sens, mais la bijection qu'il construit est en général
discontinue et non explicite --- ne pas attendre de formule
(\cref{ex:b2:structures:cbexample}).
\end{remark}

\begin{remark}[Où ce chapitre est utilisé]
Presque partout. La signature construit les déterminants
(\cref{ch:b2:linalg})~; le morphisme d'évaluation $P \mapsto P(u)$
et les idéaux principaux de $K[X]$ produisent les polynômes minimaux
et les décompositions en noyaux du~\cref{ch:b2:reduction}~; la
dénombrabilité est la scène sur laquelle se joue le~\cref{ch:b2:proba}
(les probabilités sur les espaces \omterm{def:b2:structures:countable}{dénombrables}) et la raison pour
laquelle la topologie ne cesse de produire des ensembles denses
\omterm{def:b2:structures:countable}{dénombrables} (\cref{ch:b2:metric}). La construction quotient $A/I$
est redéployée dans le volume de troisième année pour construire les
corps $K[X]/(P)$ et, à partir d'eux, la théorie de Galois~: la
propriété universelle démontrée ici y est utilisée mot pour mot.
\end{remark}

\section{Exercices}

\begin{exercise}[$\star$]\label{exo:b2:structures:1}
Lesquels des ensembles suivants sont \omterm{def:b2:structures:countable}{dénombrables}~? L'ensemble des
parties finies de $\N$~; l'ensemble de \emph{toutes} les parties de
$\N$~; $\R \setminus \Q$~; l'ensemble des polynômes à coefficients
rationnels~; l'ensemble des suites de $0$ et de $1$ nulles à partir
d'un certain rang.
\end{exercise}

\begin{exercise}[$\star$]\label{exo:b2:structures:2}
Dans $\mathfrak{S}_7$, soient $\sigma = (1\,4\,2\,6)(3\,5)$ et $\tau =
(2\,3\,7)$. Calculer $\sigma\tau$ et $\tau\sigma$ sous forme de \omterm{def:b2:structures:sn}{cycles}
disjoints, les \omterm{def:b2:structures:generated}{ordres} et signatures des quatre permutations, et
$\sigma^{2026}$.
\end{exercise}

\begin{exercise}[$\star$]\label{exo:b2:structures:3}
Construire des injections explicites montrant que $\mathcal{P}(\N)$,
$\intcc{0}{1}$ et l'ensemble $\{0,1\}^{\N}$ des suites binaires sont
deux à deux \omterm{def:b2:structures:countable}{équipotents} \emph{(développements binaires dans les deux
sens~; Cantor--Bernstein absorbe le désagrément de la double
représentation)}.
\end{exercise}

\begin{exercise}[$\star$]\label{exo:b2:structures:4}
Soit $G$ un groupe et $a, b \in G$ des éléments commutant, d'\omterm{def:b2:structures:generated}{ordres}
finis $m$ et $n$ premiers entre eux. Démontrer que
$\operatorname{ord}(ab) = mn$. Montrer par un exemple dans
$\mathfrak{S}_3$ que la commutation est essentielle.
\end{exercise}

\begin{exercise}[$\star\star$]\label{exo:b2:structures:5}
Soit $G$ un groupe fini d'\omterm{def:b2:structures:generated}{ordre} pair. Démontrer que $G$ contient un
élément d'\omterm{def:b2:structures:generated}{ordre} $2$. \emph{(Apparier chaque élément avec son inverse~;
compter ceux qui sont appariés à eux-mêmes.)}
\end{exercise}

\begin{exercise}[$\star\star$]\label{exo:b2:structures:6}
Démontrer que $A_n$ ($n \geq 3$) est engendré par les $3$-cycles.
\emph{(Un produit de deux \omterm{def:b2:structures:sn}{transpositions} est un $3$-cycle ou un
produit de deux $3$-cycles.)}
\end{exercise}

\begin{exercise}[$\star\star$]\label{exo:b2:structures:7}
Déterminer tous les morphismes de groupes~: de $(\Q, +)$ vers
$(\Z, +)$~; de $(\Z/n\Z, +)$ vers $(\Z/m\Z, +)$ \emph{(les compter~:
$\gcd(m,n)$)}~; de $(\Q, +)$ vers $(\Q_+^*, \times)$.
\end{exercise}

\begin{exercise}[$\star\star$]\label{exo:b2:structures:8}
À l'aide du théorème des restes chinois, calculer $\varphi(360)$,
trouver tous les $x$ tels que $x \equiv 3 \pmod 8$, $x \equiv 5 \pmod
9$ et $x \equiv 2 \pmod 5$, et calculer les deux derniers chiffres de
$3^{2026}$ \emph{(Euler modulo $100$~; attention~: travailler modulo
$4$ et modulo $25$)}.
\end{exercise}

\begin{exercise}[$\star\star\star$]\label{exo:b2:structures:9}
Démontrer qu'un anneau intègre fini est un corps. En déduire que
$\Z/n\Z$ est un corps si et seulement si $n$ est premier (à nouveau).
\end{exercise}

\begin{exercise}[$\star\star\star$]\label{exo:b2:structures:10}
(Un classique) Soit $K$ un corps et $G$ un sous-groupe \emph{fini} de
$(K^*, \times)$. Démontrer que $G$ est \omterm{def:b2:structures:generated}{cyclique}.
\emph{Indication~: soit $m$ l'\omterm{def:b2:structures:generated}{ordre} maximal parmi les éléments de
$G$~; montrer que l'\omterm{def:b2:structures:generated}{ordre} de tout élément divise $m$ (en utilisant
l'\cref{exo:b2:structures:4} sur des parties premières entre elles
bien choisies), de sorte que tout $G$ vérifie $x^m = 1$~; compter les
racines de $X^m - 1$.}
En particulier $(\Z/p\Z)^*$ est \omterm{def:b2:structures:generated}{cyclique}.
\end{exercise}

\begin{exercise}[$\star\star\star$]\label{exo:b2:structures:11}
Démontrer que le groupe $(\Q, +)$ n'est pas \omterm{def:b2:structures:generated}{cyclique}, et pire~: il
n'est même pas de type fini. Démontrer en revanche que tout
sous-groupe de type fini de $(\Q, +)$ est \omterm{def:b2:structures:generated}{cyclique}.
\end{exercise}

\begin{exercise}[$\star\star$]\label{exo:b2:structures:12}
(Critère de Dedekind) Démontrer que tout ensemble infini contient une
partie \omterm{def:b2:structures:countable}{dénombrable}, et en déduire qu'un ensemble $E$ est infini si et
seulement s'il est équipotent à une partie propre de lui-même.
\emph{(Pour l'implication directe, décaler une partie \omterm{def:b2:structures:countable}{dénombrable} d'un
cran~; pour la réciproque, se rappeler le principe des tiroirs.)}
\end{exercise}

\section{Problème~: le taquin}

Le taquin est un plateau $4 \times 4$ contenant quinze tuiles
coulissantes numérotées de $1$ à $15$ et une case vide~; un coup fait
glisser dans la case vide l'une des tuiles qui lui sont adjacentes.
Dans les années 1890, Sam Loyd popularisa le casse-tête en offrant
1000\$ à quiconque parviendrait à échanger les tuiles $14$ et $15$ en
ramenant toutes les autres tuiles à leur place. Personne ne les a
jamais empochés, et ce problème du week-end démontre les deux moitiés
de la raison~: la signature du~\cref{thm:b2:structures:signature}
interdit l'échange de Loyd, et --- moitié plus difficile,
constructive --- \emph{tout} ce que la signature autorise est
véritablement résoluble. L'énoncé complet est le théorème de
Johnson--Story (1879).

\begin{omfigure}
\begin{tikzpicture}[scale=0.72]
  \draw[omDef, thick] (0,0) grid (4,4);
  \draw[omDef, very thick] (0,0) rectangle (4,4);
  \foreach \k in {1,...,15} {
    \pgfmathtruncatemacro\row{int((\k-1)/4)}
    \pgfmathtruncatemacro\col{int(mod(\k-1,4))}
    \node at (\col+0.5, 3.5-\row) {$\k$};
  }
  \node[below, font=\small] at (2,-0.2) {résolu};
\end{tikzpicture}
\qquad
\begin{tikzpicture}[scale=0.72]
  \draw[omDef, thick] (0,0) grid (4,4);
  \draw[omDef, very thick] (0,0) rectangle (4,4);
  \foreach \k in {1,...,13} {
    \pgfmathtruncatemacro\row{int((\k-1)/4)}
    \pgfmathtruncatemacro\col{int(mod(\k-1,4))}
    \node at (\col+0.5, 3.5-\row) {$\k$};
  }
  \node[omThm] at (1.5,0.5) {$15$};
  \node[omThm] at (2.5,0.5) {$14$};
  \node[below, font=\small] at (2,-0.2) {la prime de Loyd};
\end{tikzpicture}

{\small La configuration résolue et la configuration $14$--$15$ de
Sam Loyd. La question à 1000\$~: des glissements légaux peuvent-ils
transformer le plateau de droite en celui de gauche~?}
\end{omfigure}

\begin{problem}[{Problème du week-end --- le théorème de résolubilité
de Johnson--Story}]
\label{pb:b2:structures:1}
Numérotons les cases de $1$ à $16$ dans l'\omterm{def:b2:structures:generated}{ordre} de lecture (de gauche
à droite, de haut en bas), de sorte que la case $k$ occupe la ligne
$i$ et la colonne $j$ avec $k = 4(i - 1) + j$. La case $16$ (en bas à
droite) est le \emph{domicile} de la case vide~; nous traitons la
case vide comme une seizième tuile, notée $b$ et identifiée au nombre
$16$. Une \emph{configuration} est une bijection $\sigma \colon
\intint1{16} \to \intint1{16}$, case $\mapsto$ contenu~; la
configuration \emph{résolue} est $\sigma = \mathrm{id}$. Partout,
$\varepsilon$ est la signature du~\cref{thm:b2:structures:signature}
et deux cases sont \emph{adjacentes} lorsqu'elles partagent une arête
du plateau.

\medskip
\noindent\textbf{Partie I --- Configurations, coups, signatures.}
\begin{enumerate}
  \item Justifier que les configurations sont exactement les éléments
        de $\mathfrak{S}_{16}$, donc qu'il y en a $16! =
        20\,922\,789\,888\,000$, et que le nombre de coups légaux
        depuis une configuration donnée est $2$, $3$ ou $4$, selon
        que la case vide se trouve dans un coin, sur un bord, ou à
        l'intérieur.
  \item Soit $\sigma$ une configuration, $p = \sigma^{-1}(16)$ la
        case de la case vide, et $c$ une case adjacente à $p$.
        Montrer que faire glisser la tuile de $c$ dans $p$ produit la
        configuration $\sigma' = \sigma \circ \tau$ avec $\tau =
        (p\ c)$, et en déduire que tout coup change la signature de
        signe~: $\varepsilon(\sigma') = -\varepsilon(\sigma)$.
  \item Colorions le plateau en damier~: $\chi(k) = (-1)^{i+j}$ pour
        la case $k$ en ligne $i$, colonne $j$. Montrer que tout coup
        change $\chi(\text{case de la case vide})$ de signe, et en
        déduire qu'une suite de coups ramenant la case vide à sa case
        de départ est de longueur paire.
  \item Montrer que
        \[
        I(\sigma) = \varepsilon(\sigma)\,
        \chi\bigl(\sigma^{-1}(16)\bigr)
        \]
        est invariant sous tout coup légal, et calculer
        $I(\mathrm{id})$.
\end{enumerate}

\medskip
\noindent\textbf{Partie II --- La prime de Loyd~: l'invariant à
l'œuvre.}
\begin{enumerate}[resume]
  \item La configuration de Loyd $\sigma_L$ coïncide avec la
        configuration résolue sauf que les cases $14$ et $15$ portent
        les tuiles $15$ et $14$. Calculer $I(\sigma_L)$ et conclure
        qu'aucune suite de coups ne relie $\sigma_L$ à la
        configuration résolue~: les 1000\$ de Loyd n'ont jamais couru
        le moindre risque.
  \item Montrer qu'exactement la moitié des configurations vérifient
        $I = +1$~: $\abs{\{\sigma : I(\sigma) = +1\}} = 16!/2$.
        \emph{(Pour une case vide fixée, apparier les configurations
        en composant avec une \omterm{def:b2:structures:sn}{transposition} fixée de deux autres
        cases.)}
  \item Montrer que tout coup est annulé par un coup légal, que
        «~$\sigma'$ est atteignable depuis $\sigma$ par des coups
        légaux~» est une relation d'équivalence, et que la classe $R$
        de la configuration résolue vérifie $R \subseteq \{I = +1\}$.
        Conclure qu'il y a au moins deux classes.
  \item Supposons la case vide au domicile~: $\sigma(16) = 16$.
        Montrer que $I(\sigma) = \varepsilon(\rho)$ où $\rho \in
        \mathfrak{S}_{15}$ est la restriction de $\sigma$ aux cases
        $1, \dots, 15$, et que toute configuration peut être amenée
        par des coups légaux à une configuration à case vide au
        domicile. Conclure~: pour prouver $R = \{I = +1\}$ il suffit
        de réaliser toute permutation \emph{paire} des quinze cases
        hors domicile par une suite de coups commençant et finissant
        avec la case vide au domicile.
\end{enumerate}

\medskip
\noindent\textbf{Partie III --- Tours de la case vide et groupe des
programmes.} Un \emph{programme} est une suite finie de coups légaux,
partant d'une configuration à case vide au domicile, dont la
configuration finale a de nouveau la case vide au domicile. Son
\emph{effet} est la permutation $\pi$ des cases définie par~: le
contenu de la case $x$ finit dans la case $\pi(x)$.
\begin{enumerate}[resume]
  \item Montrer qu'un programme exécuté depuis $\sigma$ aboutit à
        $\sigma \circ \pi^{-1}$~; qu'exécuter deux programmes l'un
        après l'autre compose leurs effets~; et que l'ensemble $H$ de
        tous les effets est un sous-groupe de $\mathfrak{S}_{15}$
        (permutations des cases $1, \dots, 15$) contenu dans le
        groupe alterné $A_{15}$.
  \item (Le tour élémentaire) Depuis la case vide au domicile, faire
        tourner la case vide autour du bloc $2 \times 2$ en bas à
        droite~: cases $16 \to 12 \to 11 \to 15 \to 16$. Montrer que
        l'effet est le $3$-cycle $(11\ 12\ 15)$, et que le tour
        inverse donne $(11\ 15\ 12)$. Les deux sont dans $H$.
  \item (Le grand tour) Vérifier que
        \[
        16 \to 15 \to 14 \to 13 \to 9 \to 5 \to 1 \to 2 \to 3
        \to 4 \to 8 \to 7 \to 6 \to 10 \to 11 \to 12 \to 16
        \]
        est un chemin fermé parcourant les seize cases (à pas
        adjacents seulement), et que son effet est le $15$-cycle
        \[
        \zeta = (15\ 12\ 11\ 10\ 6\ 7\ 8\ 4\ 3\ 2\ 1\ 5\ 9\ 13\
        14) .
        \]
        En écrivant $x_0 = 15$, $x_1 = 12$, $x_2 = 11$, \dots,
        $x_{14} = 14$ pour son \omterm{def:b2:structures:generated}{ordre} \omterm{def:b2:structures:generated}{cyclique}, vérifier que le tour
        élémentaire inverse de la question 10 est exactement $(x_0\
        x_1\ x_2)$.
  \item Démontrer la formule de conjugaison dans tout
        $\mathfrak{S}_n$~: pour une permutation $g$ et un $3$-cycle,
        \[
        g\,(a\ b\ c)\,g^{-1} = \bigl(g(a)\ g(b)\ g(c)\bigr),
        \]
        et noter que $H$, étant un groupe, est stable par
        conjugaison par ses propres éléments.
  \item En déduire que $H$ contient les quinze $3$-cycles
        \emph{consécutifs} du grand tour~:
        \[
        s_t = (x_t\ x_{t+1}\ x_{t+2}) \qquad (t \in \Z/15\Z,
        \text{ indices modulo } 15).
        \]
\end{enumerate}

\medskip
\noindent\textbf{Partie IV --- Engendrer le groupe alterné.}
\begin{enumerate}[resume]
  \item (Lemme A) Soient $s$ et $t$ des $3$-cycles dont les supports
        partagent exactement deux points, disons de supports $\{a, b,
        c\}$ et $\{b, c, d\}$. Montrer que, quitte à remplacer $s$ ou
        $t$ par son inverse (ce qui ne change rien au sous-groupe
        engendré), le produit $st$ est une double \omterm{def:b2:structures:sn}{transposition}~;
        montrer que $A_4$ ne contient aucun sous-groupe d'\omterm{def:b2:structures:generated}{ordre} $6$
        \emph{(un sous-groupe d'indice $2$ contient tous les carrés~;
        compter les $3$-cycles parmi les carrés)}~; et conclure que
        $\langle s, t\rangle$ est le groupe alterné tout entier des
        quatre lettres $\{a, b, c, d\}$.
  \item (Lemme B) Soit $X$ un ensemble de $k \geq 4$ lettres, $w
        \notin X$, et soit $G$ un sous-groupe d'un certain
        $\mathfrak{S}_n$ contenant toute permutation paire de $X$
        et un $3$-cycle $(u\ v\ w)$ avec $u, v \in X$. Montrer que
        pour tous $a, b \in X$ distincts il existe une permutation
        \emph{paire} $g$ de $X$ avec $g(u) = a$, $g(v) = b$, et en
        déduire $(a\ b\ w) \in G$.
  \item En déduire que le groupe $G$ du lemme B contient toute
        permutation paire de $X \cup \{w\}$ \emph{(utiliser
        l'\cref{exo:b2:structures:6}~: les $3$-cycles engendrent)}.
        Puis, en enchaînant les lemmes A et B le long des
        $3$-cycles consécutifs $s_0, s_1, \dots, s_{12}$ de la
        question 13, démontrer que $\langle s_0, \dots, s_{12}\rangle
        = A_{15}$.
  \item Conclure que $H = A_{15}$~: \emph{tout réarrangement pair des
        quinze tuiles est réalisable par un programme}, et $H$
        a $15!/2 = 653\,837\,184\,000$ éléments.
  \item (Le théorème de Johnson--Story, 1879) Assembler les questions
        6, 7, 8 et 17~: les configurations atteignables depuis la
        configuration résolue sont \emph{exactement} les $16!/2 =
        10\,461\,394\,944\,000$ configurations vérifiant $I = +1$~;
        et l'atteignabilité a exactement \emph{deux} classes, la
        classe de la configuration résolue et la classe du $\sigma_L$
        de Loyd. \emph{(Pour le second point, renuméroter les tuiles
        $14$ et $15$~: montrer que $\sigma \mapsto (14\ 15) \circ
        \sigma$ envoie les suites de coups sur des suites de coups et
        échange $\{I = +1\}$ avec $\{I = -1\}$.)}
\end{enumerate}

\medskip
\noindent\textbf{Partie V --- Critères, variantes, et la vue d'en
haut.}
\begin{enumerate}[resume]
  \item (Le critère pratique) Lire les quinze tuiles dans l'\omterm{def:b2:structures:generated}{ordre} de
        lecture de leurs cases, en sautant la case vide, et soit $N$
        le nombre d'inversions de cette liste~; soit $r$ la ligne de
        la case vide comptée depuis le \emph{bas}. Montrer que
        $I(\sigma) = (-1)^{N + r + 1}$, de sorte que $\sigma$ est
        résoluble si et seulement si $N + r$ est impair.
  \item (Actions de groupe) Une \emph{action} d'un groupe $G$ sur un
        ensemble $X$ est une application $G \times X \to X$, $(g, x)
        \mapsto g \cdot x$, avec $e \cdot x = x$ et $g \cdot (h \cdot
        x) = (gh) \cdot x$~; l'\emph{orbite} de $x$ est $G \cdot x$,
        et l'action est \emph{libre} lorsque $g \cdot x = x$ force $g
        = e$. Montrer que $h \cdot \sigma = \sigma \circ h^{-1}$
        définit une action libre de $H$ sur l'ensemble des
        configurations à case vide au domicile, que ses orbites sont
        exactement les classes d'atteignabilité mutuelle par
        programmes, et retrouver, à partir du décompte des orbites,
        que ces configurations se répartissent en exactement
        $15!\,/\,\abs H = 2$ classes.
  \item (L'obstruction $3 \times 3$) Montrer que le plateau $3 \times
        3$ n'admet \emph{aucun} chemin fermé visitant chaque case
        exactement une fois~: la stratégie du grand tour de la
        partie III échoue pour le taquin à huit. \emph{(Colorier en
        damier les neuf cases.)}
  \item (La réparation) Sur le plateau $3 \times 3$ avec les cases
        $1$ à $9$ dans l'\omterm{def:b2:structures:generated}{ordre} de lecture et le domicile $9$~:
        calculer les effets du tour périmétrique $9 \to 8 \to 7 \to
        4 \to 1 \to 2 \to 3 \to 6 \to 9$ (un $7$-cycle $\zeta'$
        fixant le centre $5$) et du tour de coin $9 \to 6 \to 5 \to 8
        \to 9$ (un $3$-cycle passant par le centre). En conjuguant le
        second par les puissances de $\zeta'$ et en enchaînant les
        lemmes A et B, démontrer que le groupe des programmes du
        taquin à huit est $A_8$ tout entier, donc qu'exactement
        $9!/2 = 181\,440$ des $9! = 362\,880$ configurations sont
        résolubles.
  \item (Un plateau pauvre) Soit maintenant le plateau un unique
        \omterm{def:b2:structures:sn}{cycle} de $n \geq 4$ cases portant $n - 1$ tuiles. Montrer que
        l'\omterm{def:b2:structures:generated}{ordre} \omterm{def:b2:structures:generated}{cyclique} des tuiles est invariant, que chaque classe
        d'atteignabilité a exactement $n(n - 1)$ configurations
        \emph{(les classes sont les orbites d'un \omterm{def:b2:structures:generated}{groupe cyclique}
        d'\omterm{def:b2:structures:generated}{ordre} $\operatorname{lcm}(n, n-1) = n(n-1)$)}, et qu'il y a
        $(n - 2)!$ classes --- pour $n \geq 5$ bien plus que $2$~:
        sur un plateau étroit l'invariant de parité ne capture
        presque rien, et la géométrie décide.
  \item Deux verdicts par le critère de la question 19~: le plateau
        entièrement inversé (tuiles $15, 14, \dots, 1$ dans les cases
        $1$ à $15$, case vide au domicile) et le plateau avec la case
        vide dans la case $1$ suivie des tuiles $15, 14, \dots, 1$
        dans les cases $2$ à $16$. Lequel est résoluble~?
  \item (Synthèse) La démonstration a deux piliers indépendants~: un
        \emph{invariant} ($I$, construit à partir du morphisme
        signature) montrant qu'au plus la moitié des configurations
        sont atteignables, et un théorème de \emph{génération
        explicite} ($H = A_{15}$) montrant qu'au moins la moitié le
        sont. En une phrase chacun, dire où sont intervenus~: la
        propriété de morphisme de $\varepsilon$~; le théorème de
        Lagrange~; la génération de $A_n$ par les $3$-cycles~; la
        conjugaison. Énoncer le méta-principe en une ligne.
\end{enumerate}
\end{problem}
